\documentclass[../../main.tex]{subfiles}
\begin{document}

\section{自由群}\label{section:ch11-自由群}

设 $G$ 是一个群，$\{a_\alpha\}$ 是 $G$ 中元素的一个族，其中 $\alpha\in J$。我们称元素族 $\{a_\alpha\}$ 生成 $G$，若 $G$ 的每一元素都能表示成元素族 $\{a_\alpha\}$ 的幂的乘积。如果元素族 $\{a_\alpha\}$ 是有限的，则称 $G$ 是有限生成的（finitely generated）。

\begin{definition}[自由群]\label{definition:ch11-free-group}
设 $\{a_\alpha\}$ 为 $G$ 中元素的一个族。假设每一个 $a_\alpha$ 生成 $G$ 的一个无限循环子群 $G_\alpha$。如果 $G$ 为群 $\{G_\alpha\}$ 的自由积，则称 $G$ 为自由群（free group），并且称 $\{a_\alpha\}$ 为 $G$ 的一个自由生成元组（system of free generators）。
\end{definition}

此时，$G$ 的每一个元素 $x$ 可由群族 $G_\alpha$ 中的元素给出的约化字表示，并且这种表示是唯一的。也就是说，如果 $x\neq 1$，则 $x$ 能被唯一地写成
\begin{align*}
x=(a_{\alpha_1})^{n_1}\cdots(a_{\alpha_k})^{n_k},
\end{align*}
其中，对于每一个 $i$，$\alpha_i\neq\alpha_{i+1}$ 并且 $n_i\neq 0$。（当然，$n_i$ 可能是负数。）

自由群由下面的扩展性质所刻画。

\begin{lemma}[自由群的扩展性质]\label{lemma:ch11-free-group-extension}
设 $G$ 是一个群，$\{a_\alpha\}_{\alpha\in J}$ 为 $G$ 中元素的一个族。如果 $G$ 是具有自由生成元组 $\{a_\alpha\}$ 的一个自由群，则 $G$ 满足以下条件：

对于任意给定的群 $H$ 及任何 $H$ 中的元素的一个族 $\{y_\alpha\}$，存在一个同态 $h:G\to H$ 使得 $h(a_\alpha)=y_\alpha$ 对于每一个 $\alpha$ 成立。

此外，$h$ 是唯一的。反之，如果扩展条件 $(*)$ 成立，则 $G$ 是具有自由生成元组 $\{a_\alpha\}$ 的自由群。
\end{lemma}

\begin{proof}
如果 $G$ 为自由群，则对于每一个 $\alpha$，$a_\alpha$ 生成的群 $G_\alpha$ 是一个无限循环群，从而存在一个同态 $h_\alpha:G_\alpha\to H$ 使得 $h_\alpha(a_\alpha)=y_\alpha$。然后应用引理 68.1。为了证明逆命题成立，设 $\beta$ 是一个固定的指标。根据假设，存在一个同态 $h:G\to\mathbb{Z}$ 使得 $h(a_\beta)=1$ 并且当 $\alpha\neq\beta$ 时，$h(a_\alpha)=0$。因此群 $G_\beta$ 是无限循环的。然后应用引理 68.5。
\end{proof}

上一节的结论（尤其是推论 68.6）蕴涵着以下定理。

\begin{theorem}\label{theorem:ch11-free-product-of-free-groups}
设 $G=G_1*G_2$，其中 $G_1$ 和 $G_2$ 分别是具有自由生成元组 $\{a_\alpha\}_{\alpha\in J}$ 和 $\{a_\alpha\}_{\alpha\in K}$ 的两个自由群。如果 $J$ 和 $K$ 是无交的，则 $G$ 是一个具有自由生成元组 $\{a_\alpha\}_{\alpha\in J\cup K}$ 的自由群。
\end{theorem}

\begin{definition}[元素族上的自由群]\label{definition:ch11-free-group-on-elements}
设 $\{a_\alpha\}_{\alpha\in J}$ 为任何一个加标族。设 $G_\alpha$ 表示所有形如 $a_\alpha^n$（$n\in\mathbb{Z}$）的符号的集合。通过定义
\begin{align*}
a_\alpha^n\cdot a_\alpha^m=a_\alpha^{n+m}
\end{align*}
使得 $G_\alpha$ 成为一个群，则 $a_\alpha^0$ 为 $G_\alpha$ 的单位元，并且 $a_\alpha^{-n}$ 为 $a_\alpha^n$ 的逆元。将 $a_\alpha^1$ 简记为 $a_\alpha$。我们将 $\{G_\alpha\}$ 的外自由积叫做元素族 $a_\alpha$ 上的自由群（free group on the elements $a_\alpha$）。
\end{definition}

若 $G$ 是元素族 $a_\alpha$ 上的自由群，我们通常对群 $G_\alpha$ 中的元素与其在外自由积构造中所涉及的单同态 $i_\alpha:G_\alpha\to G$ 的像不加区别。从而，将每一个 $a_\alpha$ 视为 $G$ 中的元素，并且族 $\{a_\alpha\}$ 形成 $G$ 的一个自由生成元组。

自由群与自由阿贝尔群之间有着密切的联系。为说明这一联系，我们回忆一下代数中交换子子群的概念。

\begin{definition}[交换子与交换子子群]\label{definition:ch11-commutator}
设 $G$ 是一个群。如果 $x,y\in G$，我们用 $[x,y]$ 表示 $G$ 中的元素
\begin{align*}
[x,y]=xyx^{-1}y^{-1},
\end{align*}
并且将其称为 $x$ 和 $y$ 的交换子（commutator）。由 $G$ 中的所有交换子所生成的子群称为 $G$ 的交换子子群（commutator subgroup），并记作 $[G,G]$。
\end{definition}

下面的结论是已知的，为完整起见，这里给出一个证明。

\begin{theorem}[交换子子群的性质]\label{theorem:ch11-commutator-subgroup}
对任何 $G$，交换子子群 $[G,G]$ 是 $G$ 的一个正规子群，并且商群 $G/[G,G]$ 是一个阿贝尔群。若 $h:G\to H$ 为 $G$ 到任何阿贝尔群 $H$ 的一个同态，则 $h$ 的核包含 $[G,G]$，从而 $h$ 诱导一个同态 $k:G/[G,G]\to H$。
\end{theorem}

\begin{proof}
\textbf{第一步。} 首先我们证明每一个交换子的共轭元是 $[G,G]$ 中的元素。推导如下：
\begin{align*}
g[x,y]g^{-1}&=g(xyx^{-1}y^{-1})g^{-1}\\
&=(gxyx^{-1})(1)(y^{-1}g^{-1})\\
&=(gxyx^{-1})(g^{-1}y^{-1}yg)(y^{-1}g^{-1})\\
&=((gx)y(gx)^{-1}y^{-1})(ygy^{-1}g^{-1})\\
&=[gx,y]\cdot[y,g],
\end{align*}
即这个元在 $[G,G]$ 中。

\textbf{第二步。} 我们来证明 $[G,G]$ 为 $G$ 的一个正规子群。设 $z$ 为 $[G,G]$ 中的任何一个元素，我们证明 $z$ 的共轭元 $gzg^{-1}$ 在 $[G,G]$ 中。元素 $z$ 是交换子及其逆元的乘积。由于
\begin{align*}
[x,y]^{-1}=(xyx^{-1}y^{-1})^{-1}=[y,x],
\end{align*}
所以 $z$ 实际上也等于交换子的乘积。令 $z=z_1\cdots z_n$，其中 $z_i$ 是交换子，则
\begin{align*}
gzg^{-1}=(gz_1g^{-1})(gz_2g^{-1})\cdots(gz_ng^{-1}),
\end{align*}
根据第一步可见，它是 $[G,G]$ 中元素的乘积，从而属于 $[G,G]$。

\textbf{第三步。} 我们证明 $G/[G,G]$ 是一个阿贝尔群。设 $G'=[G,G]$。我们希望证明
\begin{align*}
(aG')(bG')=(bG')(aG'),
\end{align*}
即 $abG'=baG'$。此式等价于
\begin{align*}
a^{-1}b^{-1}abG'=G',
\end{align*}
由于 $a^{-1}b^{-1}ab=[a^{-1},b^{-1}]$ 是 $G'$ 中的一个元素，所以上式成立。

\textbf{第四步。} 为了完成证明，注意：由于 $H$ 是一个阿贝尔群，所以 $h$ 将每一个交换子映为 $H$ 的单位元。于是，$h$ 的核包含 $[G,G]$，从而 $h$ 诱导出一个同态 $k$。
\end{proof}

\begin{theorem}\label{theorem:ch11-abelianization-of-free-group}
设 $G$ 是以 $a_\alpha$ 为自由生成元组的自由群，则 $G/[G,G]$ 是一个以 $[a_\alpha]$ 为基的自由阿贝尔群，这里 $[a_\alpha]$ 表示 $G/[G,G]$ 中 $a_\alpha$ 的陪集。
\end{theorem}

\begin{proof}
我们应用引理 67.7。任意给定一个阿贝尔群 $H$ 中的元素族 $\{y_\alpha\}$，存在一个同态 $h:G\to H$ 使得 $h(a_\alpha)=y_\alpha$ 对于每一个 $\alpha$ 成立。由于 $H$ 是一个阿贝尔群，$h$ 的核包含 $[G,G]$。因此 $h$ 诱导出一个同态 $k:G/[G,G]\to H$ 将 $[a_\alpha]$ 映为 $y_\alpha$。
\end{proof}

\begin{corollary}\label{corollary:ch11-free-group-rank}
如果 $G$ 是一个具有 $n$ 个自由生成元的自由群，则 $G$ 的每一个自由生成元组有 $n$ 个元素。
\end{corollary}

\begin{proof}
自由阿贝尔群 $G/[G,G]$ 的秩为 $n$。
\end{proof}

自由群有许多与自由阿贝尔群类似的性质。例如，如果 $H$ 是自由阿贝尔群 $G$ 的子群，则 $H$ 也是一个自由阿贝尔群。（对于有限秩的情形，第 67 节的习题 6 给出了一个证明的概要。一般情形的证明与此类似。）类似的结果对于自由群也成立，只是其证明相当复杂。我们将在第 14 章应用覆叠空间理论来证明它。

另一方面，自由群与自由阿贝尔群间也有许多差异。任何秩为 $n$ 的自由阿贝尔群，其子群的秩最多是 $n$。然而，类似的结论对于自由群并不成立。如果 $G$ 是一个 $n$ 个元素组成的自由生成元组的自由群，那么 $G$ 的子群的自由生成元组的基数可以大于 $n$，以至于是无限的！我们将在后面解释这种情形。

\subsection*{生成元与关系}

群论中的基本问题是对给定的两个群决定它们是否同构。对于自由阿贝尔群，这一问题已被解决：两个群是同构的当且仅当它们的基有相同的基数。类似地，两个自由群是同构的当且仅当它们的自由生成元组有相同的基数。（我们已针对有限基数的情形给出了证明。）

然而，对于一般的群，问题的答案不那么简单。只是对有限生成的阿贝尔群才有简明的答案。

如果 $G$ 是一个有限生成的阿贝尔群，那么有一个基本定理，其结论是说：$G$ 是两个子群的直和，即 $G=H\oplus T$，其中 $H$ 是秩有限的自由阿贝尔群，$T$ 是所有有限阶元素组成的子群（称 $T$ 为 $G$ 的挠子群（torsion subgroup）。）$H$ 的秩是由 $G$ 唯一确定的，因为其秩便是 $G$ 相对于它的挠子群的商群的秩。这个数值通常称为 $G$ 的 Betti 数（Betti number）。此外，子群 $T$ 自身可以表示成以素数方幂为阶的有限多个循环群的直和。这些群的阶数由 $T$ 唯一决定（从而也由 $G$ 唯一决定），它们被称为 $G$ 的初等因子（elementary divisor）。$G$ 的同构类完全由它的 Betti 数及它的初等因子决定。

如果 $G$ 不是阿贝尔群，我们还没有一个令人满意的答案，即便 $G$ 是有限生成的也是如此。关于决定 $G$ 的等价类的条件，我们所能给出的只是以下一些事实：

对于给定的 $G$，假设 $G$ 的生成元的族 $\{a_\alpha\}_{\alpha\in J}$ 已经给定。设 $F$ 是由元素 $\{a_\alpha\}$ 给出的自由群，那么，这些元素到 $G$ 的映射 $h(a_\alpha)=a_\alpha$ 可以扩充为一个满同态 $h:F\to G$。如果 $N$ 等于 $h$ 的核，则 $F/N\cong G$。因而，刻画 $G$ 的一个方法便是：给出 $G$ 的生成元族 $\{a_\alpha\}$，并且以某种方式给出子群 $N$ 的特征。$N$ 的每一个元素称为 $F$ 上的一个关系（relation），并且 $N$ 称为关系子群（relations subgroup）。我们可以通过给出 $T$ 的一个生成元的集合来描述子群 $N$。然而由于 $N$ 是 $F$ 的一个正规子群，所以还可以用一个较小的集合来描述 $N$。特别地，我们也可以用 $F$ 中元素的一个族 $\{r_\beta\}$ 来描述 $N$，只是要求这些元素及其共轭元共同生成 $N$，也就是说，$N$ 是包含这些 $r_\beta$ 的 $F$ 的最小正规子群。对于这种情形，我们也将族 $\{r_\beta\}$ 称为 $G$ 的一个关系的完备集（complete set of relation）。

由于 $N$ 的每一个元素均属于 $F$，因而它当然可以由 $G$ 的生成元 $\{a_\alpha\}$ 的幂给出的约化字唯一表示。当我们提到 $G$ 的生成元上的一个关系时，有时是指这一约化字，而不是它所表示的 $N$ 中的元素。具体含义可依照上下文来判定。

\begin{definition}[群的表示]\label{definition:ch11-group-presentation}
设 $G$ 是一个群，$G$ 的一个表示（presentation）是指 $G$ 的生成元的一个族 $\{a_\alpha\}$ 连同 $G$ 的关系的一个完备集 $\{r_\beta\}$，其中每一个 $r_\beta$ 都是集合 $\{a_\alpha\}$ 生成的自由群中的一个元素。若 $\{a_\alpha\}$ 是有限族，则 $G$ 当然是有限生成的。如果 $\{a_\alpha\}$ 和 $\{r_\beta\}$ 都是有限的，则 $G$ 称为有限表示的（finitely presented），并且这两个族一起称为 $G$ 的一个有限表示（finite presentation）。
\end{definition}

这种刻画 $G$ 的方式还远未达到令人满意的程度。$G$ 的一个表示在不区别同构的群的意义下唯一地决定了 $G$。但是，具有两个全然不同的表示的群却可能是同构的。并且，即使对于有限的情形，我们也无法确定具有不同表示的两个群是否同构。这就是群论中的“同构问题的不可解性”。

但是，我们所能做到的仅此而已！

\subsection*{习题}

\begin{exercise}\label{exercise:ch11-69-01}
若 $G=G_1*G_2$，证明
\begin{align*}
G/[G,G]\cong(G_1/[G_1,G_1])\oplus(G_2/[G_2,G_2]).
\end{align*}
[提示：对直和和自由积应用扩展条件来定义同态
\begin{align*}
G/[G,G]\rightleftarrows(G_1/[G_1,G_1])\oplus(G_2/[G_2,G_2])
\end{align*}
使得它们互逆。]
\end{exercise}

\begin{exercise}\label{exercise:ch11-69-02}
将上题的结果推广到任意自由积上。
\end{exercise}

\begin{exercise}\label{exercise:ch11-69-03}
定理 设 $G=G_1*G_2$，其中 $G_1$ 和 $G_2$ 分别是阶为 $m$ 和阶为 $n$ 的循环群，则 $m$ 和 $n$ 是由 $G$ 唯一确定的。
\begin{enumerate}[(a)]
\item 证明 $G/[G,G]$ 的阶为 $mn$。
\item 确定最大整数 $k$ 使得 $G$ 含有阶为 $k$ 的元素。（见第 68 节的习题 2。）
\item 证明定理。
\end{enumerate}
\end{exercise}

\begin{exercise}\label{exercise:ch11-69-04}
证明：若 $G=G_1\oplus G_2$，其中 $G_1$ 和 $G_2$ 分别是阶为 $m$ 和阶为 $n$ 的循环群，则一般而言，$m$ 和 $n$ 不能由 $G$ 唯一确定。[提示：若 $m$ 和 $n$ 分别为素数，证明 $G$ 是 $mn$ 阶循环群。]
\end{exercise}

\end{document}